\documentclass[12pt]{article}
\usepackage[utf8]{inputenc}
\usepackage[T1]{fontenc}
\usepackage{hyperref}
\title{Zari Language Resources for K-12 Curriculum Development}
\author{Compiled by Claude Sonnet 5.5 (claude-sonnet-5-5) for Andr\'as Kornai}
\date{}
\begin{document}
\maketitle

\section{Introduction}
Zari (ISO 639-3 zaz; Glottocode zari1242) is a dormant Chadic (Afro-Asiatic) language of Bauchi State, Nigeria. Ethnologue reports no known first-language speakers, and the ethnic population is given as about 20,700. Reported varieties are Zari, Zakshi and Boto. It should not be confused with Chaari, a separate Chadic language of the same area. The community also uses Hausa and other local languages.

\section{Primary Educational Resources}
No educational materials for Zari turned up in the searches. Any revitalization curriculum would start from scratch.

\section{Linguistic and Reference Materials}
\begin{itemize}
\item Wikipedia article on the Zari language. \url{https://en.wikipedia.org/wiki/Zari_language}
\item Endangered Languages Project entry. \url{https://endangeredlanguages.com/lang/10452}
\item A second Endangered Languages Project record that came up in the same search. \url{https://endangeredlanguages.com/lang/10485?hl=en}
\end{itemize}
These are overviews, not descriptive grammars or word lists. No substantial documentation was found.

\section{Religious and Community Resources}
None found.

\section{Digital and Audio Resources}
None found.

\section{Curriculum Development Context}
Some younger people are said to use Zari as a second language, so revitalization would rely on whatever recollection and elicitation of the language is still possible. A first step would be to search archives (for example, fieldnotes by Chadic specialists) and to consult community elders in Toro, Dass and Tafawa Balewa areas. The searches did not locate any such material.
\end{document}